% Yeban User Guide — Chapter 11: Troubleshooting (English)
% T-DOC-001 chapter-by-chapter delivery, written ahead of the CH11 task (last outline chapter).
% Sources: docs/manual chapters, AGENTS.md execution discipline, T-FIX-196 memory threshold,
% app.slint device/diagnostic entries.

\chapter{Troubleshooting}
\label{ch:troubleshooting}

\begin{quote}
人谁无过？过而能改，善莫大焉。
\end{quote}

Source: Zuo Zhuan, 2nd year of Duke Xuan (《左传·宣公二年》). Translation: “Who among us has no faults? To err and be able to change is the greatest good.”

Mistakes are not the end: follow the clues, correct the cause, and the project moves on.

This chapter is organised by symptom: what you see, then what to rule out. Each section
moves from the most likely and easiest-to-check cause towards project and device levels;
change one thing at a time and audition after each. One guiding rule throughout: Yeban
records key facts in receipts and the status bar — do not conclude from "it looks fine".

\section{No Sound}

Work through these in order, changing only one thing per step and listening after each:

\begin{enumerate}
  \item \textbf{Transport and playhead}: is playback actually running? Press Space and watch the playhead advance; Shift+Space resumes from the cursor.
  \item \textbf{Mute / solo}: are Mute or Solo engaged on the track header or in the mixer? If another track is soloed, the current one is silent by design.
  \item \textbf{Fader and pan}: is the track fader or master fader all the way down? Mix $\to$ Reset Master Fader to 0dB restores the master in one step.
  \item \textbf{Device chain}: did the instrument load? Is an effect bypassed? Bypassed devices are not the same as deleted ones.
  \item \textbf{Clip and fades}: is there actually a clip under the playhead? Was it trimmed down to silence, or are its fade-in/fade-out swallowing the whole thing?
  \item \textbf{Materials present?}: a missing embedded sample is recorded as MissingSample and the track is skipped (saving still succeeds). Check the asset pool for the audio it uses; relative SFZ paths depend on the project directory — a project opened from file uses the real directory, while an unsaved new project uses a placeholder under the output directory, which breaks samples kept only next to the project file.
  \item \textbf{Control Room}: does the master-bus meter show a level? Level present but nothing heard points to the output device; no level points to the project side.
\end{enumerate}

\screenshot{ch11-silence-checklist}{The silence checklist as presented in the help panel}

\section{Devices and Drivers}

\subsection{Input device and level}

The recording panel has input-device selection and an input test, so you can confirm a
signal arrives before actually recording. The recording path uses the project's sample
rate (48 kHz by default); after changing devices, the next recording uses the new one.

\subsection{Output and latency}

If the project-side meter shows level but you hear nothing, check the system output
device and the Control Room's master-bus reading. Yeban targets low-latency
determinism and handles internal plug-in delay compensation; for latency-sensitive
work, check the device buffer settings first rather than hunting for interface switches.

\subsection{Safe defaults}

Network-listening services (MCP HTTP) bind only the loopback address \texttt{127.0.0.1},
never \texttt{0.0.0.0}; a default build does not listen at all unless you explicitly pass
\texttt{--enable-mcp-http} or set \texttt{YEBAN\_MCP\_HTTP=1}. If MCP "won't connect" and
you are tempted to open up listening, check that runtime switch first — do not change
the bind address.

\section{A Project Will Not Open}

\begin{enumerate}
  \item \textbf{Check the file}: \texttt{.yeban} is the only project format. Saving is atomic, so a power cut mid-write usually leaves the previous complete version rather than half a file; try a backup if you have one.
  \item \textbf{Use the right entry}: File $\to$ Open / Restore Project..., or \texttt{yeban-app --open <path>}; the startup chooser's Recent Projects tab also returns you to prior work.
  \item \textbf{Are the assets alongside it?}: the project references asset-pool materials by content hash. If you copied only the \texttt{.yeban} file and lost referenced materials, the project opens but those tracks are silent — that is not "won't open".
  \item \textbf{Export a diagnostic bundle}: Diagnostics $\to$ Export Project Diagnostic Bundle produces a troubleshooting bundle (not artwork audio). Attaching it is far more useful than a vague description when asking for help.
  \item \textbf{Record errors verbatim}: if the app reports an error code or a missing list, note the original text; do not substitute "it looks like the same problem as last time" for a measurement.
\end{enumerate}

\section{Performance and Memory Thresholds}

\subsection{Know the threshold}

The app test target is heavy: measured on a 4 GB machine, \texttt{cargo test -p yeban-app}
is killed by the kernel (rustc SIGKILL) at roughly 35 minutes, yielding no results; moving
the same kind of assertions into a lighter crate (\texttt{yeban-model}) compiled in about
1m39s and ran in about 1s. The practical threshold for app-side verification is at least
8 GB; below it, app-side work can only be static changes — do not claim tests were run.

\subsection{A minimum verification set when resources are tight}

When memory is insufficient, use the smallest set you can actually complete:
\texttt{cargo fmt --check}, \texttt{cargo check --all-targets} and clippy for the target
crate. Record what could not run honestly as \texttt{not\_run} with the reason, and let CI
rerun it on a larger machine. For building, use
\texttt{cargo build --profile release-small -p yeban-app}; on a headless machine add
\texttt{--headless}.

\subsection{Only the CI verdict counts}

A local green is only a reference; the final verdict belongs to CI. An unread verdict
counts as pending, never as "passed". Also avoid treating an overall "success" as proof
that the target assertions actually ran — confirm the job executed (non-zero duration)
and its log contains test-result lines.

\section{Help Topic Mapping}

\begin{table}[htbp]
\centering
\caption{Sections of this chapter $\to$ in-app help anchors}
\label{tab:ch11-help}
\begin{tabular}{ll}
\hline
Section & Help topic \\
\hline
No sound & \texttt{troubleshoot-silence} \\
Devices and drivers & \texttt{troubleshoot-devices} \\
A project will not open & \texttt{troubleshoot-project-open} \\
Performance and memory thresholds & \texttt{troubleshoot-performance} \\
Minimum verification and CI verdict & \texttt{troubleshoot-verification} \\
\hline
\end{tabular}
\end{table}

\paragraph{What you will be able to do} Work through silence, device and project-open failures symptom by symptom, know the 8 GB threshold and the minimum verification set, and understand that the CI verdict is final.

\paragraph{Hands-on exercise} Reproduce "no sound": mute a track then unmute it, bypass then restore an effect, auditioning after each step and naming which step restored sound; then export a diagnostic bundle and confirm it is a diagnostic file, not the piece; finally run the minimum verification set in a low-memory environment and record the unrunnable items honestly as \texttt{not\_run}.
